\chapter{Lipoma}

\section*{Definition}
A lipoma is a benign, slow-growing mesenchymal tumor composed of mature adipocytes, typically presenting as a soft, mobile, painless subcutaneous nodule. They are the most common soft-tissue tumor, occurring in approximately 1\% of the adult population.

\section*{Pathophysiology}
Lipomas arise from primitive adipose cells during postnatal life, with clonal proliferations of mature adipocytes driven by chromosomal translocations involving the HMGA2 gene (12q13--15) in 50--60\% of cases. These tumors are surrounded by a thin fibrous capsule, contain no atypical cells, and lack malignant potential. Growth follows a slow, self-limited course, often reaching a stable size of 2--10 cm.

\section*{Etiology}
\begin{itemize}
  \item \textbf{Sporadic:} Majority are idiopathic with no identifiable trigger
  \item \textbf{Genetic:} Familial multiple lipomatosis (autosomal dominant), Gardner syndrome, Proteus syndrome, Cowden syndrome, Dercum disease (adiposis dolorosa)
  \item \textbf{Risk factors:} Middle age (40--60 years), obesity, diabetes mellitus, alcohol abuse (some studies)
  \item \textbf{Trauma:} Rare association with prior blunt trauma (post-traumatic lipoma)
\end{itemize}

\section*{Indications for Evaluation}
Seek care if:
\begin{itemize}
  \item Rapid growth, pain, or tenderness at the site
  \item Size \ensuremath{>}5 cm, deep location, or fixation to underlying tissues (suspect liposarcoma)
  \item Functional impairment, compression of neurovascular structures, or cosmetic concern
  \item New or enlarging lipoma in a patient with known familial cancer syndrome
\end{itemize}

\section*{Related Symptoms}
\begin{itemize}
  \item Liposarcoma, angiolipoma, neurofibroma, epidermal inclusion cyst, sebaceous cyst, hibernoma
\end{itemize}
\textbf{Note:} The probability of liposarcoma increases with size \ensuremath{>}5 cm, deep (subfascial) location, rapid growth, and older age; any suspicious lesion warrants imaging and/or biopsy.

\section*{Self-Care / Home Management}
\begin{itemize}
  \item No self-care is needed for asymptomatic lipomas; observation is appropriate
  \item Weight loss does not reduce lipoma size since lipomas are not metabolically responsive
  \item Avoid trauma or manipulation as it may cause inflammation
\end{itemize}

\section*{Diagnostic Approach}
\begin{itemize}
  \item Clinical diagnosis: Soft, rubbery, mobile, non-tender mass with smooth borders and normal overlying skin
  \item Imaging: Ultrasound (well-defined, homogeneous, compressible, avascular), CT or MRI for deep or suspicious lesions
  \item Biopsy/excision for histologic confirmation if any atypical features are present
\end{itemize}

\section*{Treatment / Management Overview}
\begin{itemize}
  \item Observation: Asymptomatic and cosmetically acceptable lipomas require no intervention
  \item Excision: Simple surgical enucleation with capsule intact; recurrence rate <2\% for complete excision
  \item Liposuction-assisted removal for large lipomas or cosmetic body contouring
  \item Steroid injection (triamcinolone) may reduce size in small inflammatory or angiolipomas
\end{itemize}

\clearpage
